\documentclass[11pt,a4paper]{article}
\usepackage[margin=27mm]{geometry}
\usepackage{amsmath,amssymb,amsthm,hyperref}
\newtheorem{proposition}{Proposition}
\newtheorem{theorem}[proposition]{Theorem}
\newcommand{\gen}[1]{\langle #1\rangle}
\title{Kernels and intersections in quotients of hypertopes}
\author{Human Compression project}
\date{}
\begin{document}
\maketitle
\begin{abstract}
Injectivity on the maximal parabolic subgroups does not by itself preserve
the intersection property under a quotient. We give an affine example in
which every maximal parabolic embeds and the quotient geometry is flag
transitive, but has only one chamber. We then state and prove the precise
additional condition on products of parabolic subgroups. This corrects
the quotient assertion in Theorem 3.2 of the fixed manuscript
\emph{An exploration of locally spherical regular hypertopes}.
It does not dispose of that manuscript's classification results.
\end{abstract}

\section{The direction of a quotient criterion}

Let $I$ be a finite set, and let $U=\gen{s_i:i\in I}$ be generated by
involutions. Write $U_A=\gen{s_i:i\in A}$ for $A\subseteq I$ and
$U_i=U_{I\setminus\{i\}}$. The marking is a C-group marking when
\[
 U_A\cap U_B=U_{A\cap B}\qquad(A,B\subseteq I).
\]
All subgroups in this note refer to the specified marking.

The source's Theorem 3.1 is a valid upward criterion: if a marked
epimorphism $G\longrightarrow H$ is injective on each maximal parabolic
of $G$, and $H$ is a C-group, then $G$ is a C-group. Indeed, its proper
parabolics inherit the property from their isomorphic images in $H$.
An element of $G_i\cap G_j$ maps into $H_{I\setminus\{i,j\}}$;
lift it within $G_{I\setminus\{i,j\}}$ and use injectivity in $G_i$.
The maximal-parabolic criterion proved below completes the argument.

The proof of Theorem 3.2 instead uses a map $U/N\longrightarrow U$
which takes $s_iN$ to $s_i$. Local injectivity does not make this map
well defined: an additional relation in the quotient need not hold in
$U$. The issue is not confined to that proof. The following example
satisfies the stated hypotheses and contradicts its conclusion.

\section{An affine example}

Equip $\mathbb R^2$, in coordinates $(a,b)$, with the inner product
whose matrix is $\left(\begin{smallmatrix}2&-1\\-1&2\end{smallmatrix}\right)$.
Use homogeneous coordinates for the affine maps
\[
 s_0=\begin{pmatrix}0&-1&1\\-1&0&1\\0&0&1\end{pmatrix},\quad
 s_1=\begin{pmatrix}-1&1&0\\0&1&0\\0&0&1\end{pmatrix},\quad
 s_2=\begin{pmatrix}1&0&0\\1&-1&0\\0&0&1\end{pmatrix}.
\]
They are reflections in $a+b=1$, $2a-b=0$, and $-a+2b=0$,
respectively. Put $U=\gen{s_0,s_1,s_2}$ and $W=\gen{s_1,s_2}$.
The latter is the group of six linear symmetries generated by the
last two reflections. If $t_v$ denotes translation by $v$, then
\[
 s_\theta=s_1s_2s_1=\begin{pmatrix}0&-1&0\\-1&0&0\\0&0&1\end{pmatrix},
 \qquad s_0=t_{(1,1)}s_\theta.
\]
Conjugating $t_{(1,1)}$ by $s_1$ gives $t_{(0,1)}$, and their quotient
gives $t_{(1,0)}$. Every generator has integral translation part and
linear part in $W$. Consequently
\begin{equation}\label{semidirect}
 U=\mathbb Z^2\rtimes W,
 \qquad N=\{t_v:v\in\mathbb Z^2\}\triangleleft U,
 \qquad U/N\cong W\cong S_3.
\end{equation}

We verify that $U$ is the full type-preserving group of a regular
hypertope, so no converse about arbitrary marked groups is being
assumed. The lines
\[
 2a-b\in\mathbb Z,\qquad -a+2b\in\mathbb Z,\qquad a+b\in\mathbb Z
\]
tile the plane by equilateral triangles. A fundamental triangle is
\[
 \Delta=\{2a-b\geq0,\ -a+2b\geq0,\ a+b\leq1\},
\]
with vertices $0$, $\lambda_1=(2/3,1/3)$, and
$\lambda_2=(1/3,2/3)$. The vertex set is the disjoint union of
$\mathbb Z^2$, $\lambda_1+\mathbb Z^2$, and
$\lambda_2+\mathbb Z^2$. Give these three sets different types.
Each generator preserves each set, as follows directly from the
displayed matrices. Declare vertices incident when they coincide or
are joined by a tiling edge.

Each edge belongs to two triangles; each vertex link is a hexagon;
every three pairwise adjacent vertices form a tiling triangle.
It follows that this incidence structure is a geometry, is thin,
and is residually connected. Here thinness concerns extensions of
an edge by a vertex of the third type, and residual connectivity
concerns the whole tiling graph and its hexagonal vertex links.

The group $U$ is transitive on triangles. To see this, cross a path
of adjacent triangles starting at $\Delta$. Once the current triangle
is $u\Delta$, reflection in one of its sides is $us_iu^{-1}$;
the next triangle is therefore $us_i\Delta$.
The adjacency graph is connected, so induction reaches every triangle.
The action is free on triangles. An element stabilizing $\Delta$
preserves types and hence fixes each of its three noncollinear
vertices, which forces an affine map to be the identity.
Finally, a type-preserving automorphism fixing one chamber fixes each
neighboring chamber, since an edge has exactly two extensions.
Propagation across the connected chamber graph fixes the whole
geometry. Thus $U$ is its full type-preserving automorphism group,
acting regularly on chambers. Denote this regular hypertope by
$\mathcal U$.

\begin{proposition}\label{failure}
For each $i\in\{0,1,2\}$, the subgroup $U_i$ has order six and
$N\cap U_i=1$. Nevertheless the coset geometry with group $U/N$
and maximal marked parabolics $\pi(U_i)$ has one face of each type.
It is flag transitive and is not a hypertope.
\end{proposition}
\begin{proof}
Any pair of the three reflection matrices consists of involutions
whose product has order three. A group generated by two such
involutions is a quotient of the dihedral group of order six:
every word reduces to $(s_is_j)^k$ or $(s_is_j)^ks_i$, $0\leq k<3$.
The linear parts of each pair are distinct transpositions generating
the six-element group $W$. Hence each pair generates exactly six
elements and the linear-part map is injective on that subgroup.
This proves $N\cap U_i=1$ and $\pi(U_i)=W$ for all $i$.

Each coset set $W/\pi(U_i)$ is therefore a singleton. With types
kept distinct, all three faces are incident and there is exactly
one chamber. Every flag extends to it and $W$ is transitive on flags
of each prescribed type. Every residue of rank one has one face,
so the geometry is not thin. The action of $W$ has kernel $W$;
faithfulness was not a hypothesis of the source theorem.
The orbit quotient $\mathcal U/N$ is exactly the same geometry,
since translations are transitive on vertices of each type.
\end{proof}

\section{The additional kernel condition}

\begin{proposition}\label{maximal}
Suppose $G=\gen{r_i:i\in I}$, with $|I|\geq2$, and all maximal
parabolics $G_i$ are C-groups on their indicated generators.
Then $G$ is a C-group if and only if
$G_i\cap G_j=G_{I\setminus\{i,j\}}$ for all distinct $i,j$.
\end{proposition}
\begin{proof}
Necessity follows from the definition. For sufficiency, let
$x\in G_A\cap G_B$. If either index set is $I$, the conclusion is
immediate. Otherwise choose $i\notin A$ and $j\notin B$.
If $i=j$, the desired intersection follows inside the C-group $G_i$.
If $i\ne j$, the assumed maximal intersection places $x$ in
$G_{I\setminus\{i,j\}}$. Intersecting with $G_A$ inside $G_i$
places $x$ in $G_{A\setminus\{j\}}$. Both this subgroup and $G_B$
lie in $G_j$, whose intersection property now gives
$x\in G_{A\cap B}$. The reverse inclusion always holds.
\end{proof}

\begin{theorem}\label{kernel}
Let $U=\gen{s_i:i\in I}$ be a C-group with $|I|\geq2$,
let $N\triangleleft U$, and assume $N\cap U_i=1$ for every $i$.
The induced marking on $U/N$ is a C-group marking if and only if
\begin{equation}\label{condition}
 N\cap U_iU_j=\{1\}\qquad(i\ne j).
\end{equation}
In particular, if \eqref{condition} holds and the quotient coset
geometry is flag transitive, it is a regular hypertope.
\end{theorem}
\begin{proof}
Write $\pi:U\longrightarrow U/N$ and $H_i=\pi(U_i)$.
The local injectivity assumption identifies $H_i$, with its marking,
with the C-group $U_i$.

Suppose \eqref{condition} holds, and let
$x\in H_i\cap H_j$. Choose $a\in U_i$ and $b\in U_j$ with
$\pi(a)=x=\pi(b)$. Then $ab^{-1}\in N\cap U_iU_j$, so $a=b$.
As $U$ is a C-group, this common element lies in
$U_i\cap U_j=U_{I\setminus\{i,j\}}$. Thus
$H_i\cap H_j=\pi(U_{I\setminus\{i,j\}})$.
Proposition~\ref{maximal} gives the full intersection property.

Conversely suppose $U/N$ is a C-group. Write an element of
$N\cap U_iU_j$ as $ab^{-1}$ with $a\in U_i$, $b\in U_j$.
The common image of $a$ and $b$ lies in
$H_i\cap H_j=\pi(U_{I\setminus\{i,j\}})$. Choose a preimage
$c$ in that smaller subgroup. Injectivity on $U_i$ gives $a=c$,
and injectivity on $U_j$ gives $b=c$. Hence $ab^{-1}=1$.

For the final assertion, apply the C-group reconstruction theorem
under its flag-transitivity hypothesis. This is Theorem 2.3 of
the fixed source, whose complete proof is acquired from the project's
reviewed manuscript on residue stabilizers and distinguished generators.
\end{proof}

The missing condition can be seen explicitly in our example:
$s_\theta\in U_0$, $s_0\in U_1$, but
$s_\theta s_0=t_{(-1,-1)}$ is a nonidentity element of
$N\cap U_0U_1$. Thus the repair excludes the example for the
precise reason isolated in the proof.

\paragraph{Scope and verification.}
The fixed source is a sixteen-page revised author manuscript.
The affected assertion and its proof are on page 6. Proposition~\ref{maximal} is its Proposition 2.1,
attributed there to earlier work, and is reproved here for completeness.
The affine group and its triangular complex are standard; the
explicit argument above supplies all properties needed by the example.
The kernel test is an elementary quotient-intersection criterion;
no claim of historical priority is made. The accompanying program
checks the displayed integer matrices, all three six-element
parabolics, their exact linear images, and the nonzero kernel witness.
The claims concerning the infinite tiling use the written proof.
This note corrects one auxiliary theorem. It is neither a solution
of an open conjecture nor a complete treatment of the source's
spherical, Euclidean, and hyperbolic classifications.

\begin{thebibliography}{9}
\bibitem{source}
M. E. Fernandes, D. Leemans, and A. Ivi\'c Weiss,
\emph{An exploration of locally spherical regular hypertopes}.
Fixed sixteen-page revised author manuscript, catalogue L033,
especially Propositions 2.1--2.2 and Theorems 2.3, 3.1--3.2.
\bibitem{residues}
Human Compression project, \emph{Residue stabilizers and distinguished generators}.
Complete reviewed L052 general component; only the reconstruction
theorem with its stated flag-transitivity hypothesis is used here.
\end{thebibliography}
\end{document}
